\subsection{Antecedentes del proyecto de innovacion / Mejora / Creatividad}


\subsubsection{Antecedentes Nacionales}

% Tesis nacional 1
En la tesis de pregrado publicado en la Universidad Tecnologica del Perú, presentado por \textcite{aguilar2022desarrollo}, presenta como titulo de investigación ``Desarrollo de un sistema web para optimizar el proceso de entrega de boletas de pago de planillas, basado en el marco Scrum, en Sistemas y Fluidos SAC, 2021'', donde menciona que tiene como problema por planillas en la empresa de Sistemas y Fluidos S.A.C. la entrega de boletas se realizaba de manera manual. El área de recursos humanos exportaba la planilla de pagos a Excel, generaba las boletas de pago con una plantilla de Word e imprimía 2 copias de cada documento. Este proceso podía tardar hasta dos dias y estaba expuesto a producir errores. Además, gran parte de la planilla de colaboradores trabajaba en campo y no podian recogerlas en la oficina central por falta de tiempo y a la distancia, provocando reclamos de los colaboradores por la incertidumbre sobre sus deducciones y otros conceptos consignados que se encuentran en las boletas. Como objetivo decidieron desarrollar un sistema web, basado en el marco de trabajo SCRUM, para optimizar las entregas de las boletas de pago mediante la validación de la planilla, la emisión de las boletas y la generación de reportes. Para la metodología de desarrollo, se compararon Scrum, XP y RUP mediante criterios de ciclo de vida, planificación, calidad, procesos y herramientas. Al final del analisis, se selecciono SCRUM por obtener la mayor calificación con un 94.6 \%. Tras implementar el sistema web, el tiempo de preparación para cada boleta se redujo de 5,5-6 minutos a 2,25-2,58 minutos. Las boletas con información válida aumentaron de 75-93,33 \% a 98,33-100 \%, y la entrega de boletas alcanzo un 100 \% de los 60 colaboradores observados. La implementación del sistema web optimizó la entrega de boletas, mejoró la validación de la información, redujo costos y tiempo, y permitió dar un seguimiento a la recepción de los documentos por parte de los colaboradores.

% Tesis nacional 2

En la tesis de pregrado de la Universidad San Ignacio de Loyola,  sustentato por \textcite{quispe2025implementacion}, titulado ``Implementación de mejoras tecnológicas en la gestión de planillas de remuneraciones y su control presupuestario en la empresa Ascensores S.A. - Lima, Perú 2024'', se menciona que en la empresa Ascensores S.A., se detectaron deficiencias tecnológicas en la gestión de planillas y el control presupuestario. El análisis lo relacionó con tareas manuales, la falta de integración entre las áreas de Recursos Humanos y Contabilidad, y un seguimiento insuficiente de los gastos en las remuneraciones. Como objetivo decidieron implementar mejoras tecnológicas para optimizar la gestión de planillas y fortalecer su control presupuestario mediante el desarrollo de un sistema automatizado y conectado entre las áreas de Recursos Humanos y Contabilidad, el establecimiento de un modelo de análisis mensual de la planillas y el desarrollo de un plan de capacitación. El proyecto empleo una metodología con un enfoque cuantitativo y cualitativo en base al análisis de los datos históricos de las planillas y entrevistas al personal administrativo. Como resultado, obtuvieron un aumento en la precisión de la elaboración de planillas del 94.5 \% al 96 \% y un cumplimiento de los plazos de pago del 97.2 \% a 98 \% entre los años 2023 y 2024. Tambien se reportó un ahorro calculado de S/. 145,359 en el año 2023 y una proyección acumulada de S/. 499,592 para el año 2026. La autora concluye que la implementación de estas soluciones redujo los errores, mejoró la eficiencia administrativa y fortaleció la toma de decisiones. Además, considera que la automatización e integración de los sistemas fueron herramientas claves para la sostenibilidad y competitividad de la empresa.

% Tesis Nacional 3

En esta tesis de pregrado de la Universidad San Ignacio de Loyola, sustentado por \textcite{pedroza2025propuesta}, presentado con el titulo de ``Propuesta de optimización del sistema de información para la gestión de personas en PROMPERÚ'', donde se menciona que en la entidad conocida como PROMPERÚ, presenta deficiencias en los procesos de recursos humanos a causa de que muchos de ellos son rudimentarios o no están digitalizados; los datos necesarios están dispersos y, sumado al factor del error humano, esto aumenta la carga laboral del departamento de recursos humanos, tomando un tiempo mayor del esperado. Se propone la implementación de un sistema para automatizar y digitalizar los procedimientos manuales anteriormente mencionados, con el fin de reducir los errores humanos y lograr tiempos de revisión mucho menores en comparación con el uso de procedimientos rudimentarios. Se utilizó la metodología BPM (Business Process Management), empleada en la gestión de procesos, la cual resulta ideal para este proyecto al centrarse en la optimización y el monitoreo. Se logró la implementación del sistema para la optimización de la gestión de personas dentro de la empresa PROMPERÚ; se redujeron las fuentes de datos de 14 a 1 sola fuente unificada, lo que dio como resultado una reducción del 80 \% en el sobretiempo de los empleados dedicado a los procesos administrativos. La implementación de un sistema de optimización de procesos reduce significativamente la demora al ejecutar tareas que anteriormente se realizaban de forma manual.

% Tesis nacional 4

En este proyecto de la Facultad de Ingenieria de Sistemas e Informática de la Escuela Profesional de Ingenieria de Sistemas e Informática, sustentado por \textcite{fernandez2023implementacion}, presentado con el titulo de ``Implementación de un sistema web de remuneraciones para optimizar la elaboración de planillas de la empresa BIP BIP, 2023'', se menciona que dentro de la empresa llamada BIP BIP, en el area de Recursos Humanos se detectaron demoras significativas en procesos como generacion de planillas, tomas de asistencias, flata de transparencia y dificultades al realizaar el cumplimiento normativo. Cuyo objetivo cual es el uso de un sistema informatico para la optimizacion de procesos , que contribuya a la mejora de gestion por parte del area de recursos humanos y la del personal de la entidad. Se empleó la metodología RUP (Proceso Unificado de Rational), para identificar los puntos críticos y la posterior creación de soluciones adecuadas siguiendo los requerimientos identificados bajo la arquitectura definida del proyecto. Se redujo significativamente el tiempo de elaboracion de procedimientos tales como, planillas, toma de asistencias, generacion de boletas mediante la implementacion del sisitema web. La implementación del sistema informatico web, logro reducir el tiempo promedio de demora al realizar procedimientos como, toma de asistencia, elaboracion de planillas y boletas, logrnado los objetivos del projecto.

% Tesis nacional 5

En esta tesis de pregrado de la Universidad Nacional Toribio Rodríguez de Mendoza de Amazonas, sustentado por \textcite{suarez2025sistema}, titulado ``Sistema web para optimizar el proceso de gestión de planillas del supermercado RB de Bagua'', se menciona que en el Supermercado RB de Bagua, se presentaban deficiencias relacionadas con el registro manual de asistencias, duplicados de datos y errores de cálculos al realizar pagos a los empleados. El objetivo es implementar un sistema web que permita la automatización de las tareas administrativas, mejorar la precisión de los cálculos para evitar el error humano y garantizar el seguimiento laboral. Se aplicó la metodología SCRUM debido a la facilidad para coordinar trabajos múltiples dentro del propio equipo, así como la independencia que da dentro del proyecto. Se mostraron grandes mejoras en todas las áreas evaluadas donde se presentaban las problemáticas, teniendo una opinión positiva por parte de los empleados y de la entidad donde se aplicó el sistema web. La incorporación del sistema web redujo drásticamente el error humano y agilizó procesos, logrando una mayor satisfacción del personal al usarlo.

\subsubsection{Antecedentes Internacionales}
%antecedente nro 1
En este proyecto de grado de especialización de la Facultad de Ingenierías del Instituto Tecnológico Metropolitano (ITM), presentado por \textcite{arias2026automatizacion} titulado ``Automatización del Sistema de Gestión de Nómina de Contratistas del Instituto Tecnológico Metropolitano``, cual se realizo un análisis se identificó que en el Instituto Tecnológico Metropolitano (ITM), gestionaba la nómina mensual de 700 contratistas mediante correos electrónicos y hojas de cálculo, sin un sistema que centralizara los datos, validara los documentos recibidos y registrara las acciones realizadas.	Se definió el objetivo de implementar un sistema de gestión de nóminas de contratistas que abarque desde la carga documentos por cada dependencia hasta la generación automática de los archivos necesarios para el pago.	Este trabajo aplicó una metodología de desarrollo iterativo e incremental. Recopilaron los requisitos mediante una entrevista a la liquidadora de nómina y análisis de documentos del proceso existentes; después se diseño la arquitectura con el método ADD, se desarrolló un prototipo y se validaron sus características de calidad.	Como resultado, se documentaron once requisitos funcionales de prioridad alta y se diseño la arquitectura. También se desarrollo un prototipo de interfaz con funciones de carga y validación de documento. El prototipo trabajo con datos simulados y no estaba integrado con el backend; por ello, el trabajo no presenta cifras.	Los autores concluyeron que la arquitectura propuesta y el prototipo muestran que resuelven las deficiencias encontradas en el proceso actual como las inconsistencias, rechazos bancarios, entre otros. Asimismo, el sistema demuestra la viabilidad técnica y arquitectónica para el contexto de la empresa.
%antecedente nro 2

En este proyecto de tesis de grado, publicado en la Universidad Pública de El Alto y presentado por \cite{zeballos2022sistema}, titulado ``Sistema De Información Web Para El Control De Acceso Y Asistencia Para El Proceso De Planillas Mediante RFID'', se identificó que la empresa consultora carecía de un sistema automatizado para el procesamiento de planillas y el registro de personal. La empresa dependía de plantillas básicas en Excel con macros en VBA que generaban retrasos considerables, aumento en los costos de operación, riesgos de pérdida de información por archivos corruptos y reclamos del personal debido a discrepancias en el cómputo de tardanzas y ausencias. El objetivo fue desarrollar un sistema de información web para el control de asistencias y accesos mediante tecnología RFID que permita automatizar la generación de planillas de sueldo y la emisión de boletas de pago de manera eficiente y precisa. Se optó por un enfoque ágil combinando la metodología Kanban con la ingeniería web basada en diagramas UML para el modelado de navegación, contenido y presentación. Para la arquitectura técnica se definió un API RESTful, empleando el framework Laravel (PHP) para el backend, un frontend SPA reactivo desarrollado con Vue.js y Quasar, y MariaDB como base de datos. La implementación de un motor contable automatizó el cálculo de salarios netos integrando días/horas trabajadas, recargo por horas extraordinarias, bonos y deducciones de ley. La evaluación bajo normas ISO/IEC 25000 certificó una adecuación funcional del 84.44\%, usabilidad del 94\% y una calidad global del sistema del 88.8\%, reduciendo el margen de error en la generación de boletas de pago en PDF a cero. La adopción de una arquitectura web desacoplada (API REST con backend y frontend separados) articulada con bases de datos relacionales permite centralizar de forma segura la trazabilidad de incidencias de asistencia y sincronizarlas directamente con el cálculo de la planilla, garantizando transparencia contable y cumplimiento estricto de las normativas laborales vigentes en el país.

%antecedente nro 3

En este proyecto de integración curricular, presentado en la Escuela Politécnica Nacional por \cite{masapanta2025desarrollo}, titulado ``Desarrollo de un módulo integral para la gestión de nómina y talento humano en la Cooperativa de Ahorro y Crédito Politécnica Ltda.'', se identificó una necesidad en el área de recursos humanos de la institución relacionada con la gestión del personal, específicamente al administrar los datos requeridos para el cálculo de nóminas, así como la gestión de vacaciones y del expediente de los empleados. El objetivo principal de este proyecto es la integración de un módulo para la administración del talento humano y nómina, que agilice la gestión de los datos de los integrantes de la empresa y automatice la creación, gestión y descarga de los pagos de estos. Se optó por la metodología Scrum, la cual permite la identificación, creación y distribución de los requerimientos según el rol y prioridad en plazos fijos, bajo un entorno colaborativo que garantiza un avance continuo. Se logró la implementación del módulo de manera exitosa, cubriendo las necesidades del área de recursos humanos en la administración de nóminas, vacaciones y expedientes. Tras pruebas exhaustivas en el entorno laboral, se confirmó la integración correcta del módulo, aplicando las reglas de negocio y cumpliendo con los requerimientos identificados inicialmente.

%antecedente nro 4

En este trabajo de grado, publicado en la Universidad El Bosque y presentado por \cite{beltran2021automatizacion}, titulado ``Automatización de procesos encargados de gestionar la operación del personal, los clientes y cálculo de nómina de la empresa E.D.S San Antonio Cogua mediante su análisis procesal'', se presenta una solución a la entidad E.D.S. San Antonio frente a la problemática de sus procesos de administración interna, los cuales se realizaban de manera manual, causando retrasos, dependencia de roles específicos, errores humanos y duplicidad de actividades. El objetivo principal fue la creación de un aplicativo web que automatice los procedimientos de administración interna, la gestión de datos de empleados, los pagos nominales y la atención a clientes. Para su desarrollo se aplicó el modelo en Cascada para la estructuración secuencial del software, la metodología BPM (Business Process Management) para el análisis de requerimientos y la planificación de actualizaciones, y la Revisión Sistemática de la Literatura para la selección de tecnologías de código abierto (\textit{open source}). Con la implementación del sitio web se redujeron en más del 50\% los tiempos de ejecución de los procesos administrativos y se evidenció un incremento del 79\% en la productividad, cumpliendo con éxito los requerimientos e incrementando la eficiencia operativa de la empresa.

% atenecedente nro 5

En este proyecto final, publicado en la Universidad Tecnológica Nacional (Facultad Regional Rosario) y presentado por \cite{basualdo2026sistema}, titulado ``Sistema de Gestión de Recursos Humanos'', se identificaron ineficiencias estructurales en las organizaciones derivadas de la gestión manual y dispersa de los procesos de recursos humanos. Entre los principales problemas destacaban la ausencia de trazabilidad en las liquidaciones salariales, el riesgo de errores aritméticos en la asignación de deducciones e incentivos, y las dificultades operativas para sincronizar la asistencia laboral con el cómputo final de la nómina y los perfiles de los colaboradores. El objetivo principal fue diseñar, desarrollar y evaluar la viabilidad integral de un Sistema de Gestión de Recursos Humanos bajo el modelo SaaS (\textit{Software as a Service}) que automatice el ciclo de vida del empleado, integre el control de tiempos y asistencias con el motor de liquidación salarial, y provea mecanismos de auditoría continua y analítica de datos mediante un panel gerencial (\textit{dashboard}). Para la gestión del desarrollo se utilizó la metodología ágil Kanban, dimensionando el esfuerzo técnico mediante Puntos de Casos de Uso (114 PCU calculados, equivalentes a 1,368 horas de desarrollo) y la formalización de 44 Casos de Uso detallados bajo estándares UML. La arquitectura técnica se estructuró con un backend en Python/Django, un frontend responsivo mediante Bootstrap y JavaScript, y una base de datos relacional (PostgreSQL/MySQL), incorporando control de accesos basado en roles (RBAC) y un registro estricto de eventos (\textit{audit logs}) alineado con la norma ISO/IEC 27001. La evaluación financiera demostró la rentabilidad del software con un Valor Actual Neto (VAN) positivo de \$4,103.23 USD y una Tasa Interna de Retorno (TIR) anual del 21.11\%, confirmando la sostenibilidad del modelo operativo. Asimismo, la plataforma logra integrar legajos digitales, asignación cruzada de beneficios, liquidación masiva de nóminas y paneles analíticos en tiempo real, garantizando un entorno informático escalable, seguro y alineado con los estándares de ciberseguridad.


\subsubsection{Artículo Científico}

% Articulo 1 RF 13

En este artículo de investigación publicado por \textcite{ahmed2023web}, titulado ``Web-based payroll management system: design, implementation, and evaluation'', 
donde mencionan que tiene como problemática que la gestión contable de salarios en múltiples organizaciones y facultades se realiza de forma manual y en soporte físico, 
lo que genera retrasos considerables en el cálculo de deducciones y asignaciones, vulnerabilidad a errores aritméticos humanos, dificultades para consultar el histórico de períodos anteriores, 
retrasos en la generación de boletas y baja confiabilidad en los reportes financieros consolidados. 
Su objetivo es diseñar, desarrollar y evaluar un sistema integral de gestión de nóminas basado en web (WPMS) que automatice el cálculo de compensaciones mensuales y anuales, 
centralice el archivo de datos salariales en un repositorio histórico, permita el control de accesos segregado por roles y agilice la emisión y fiscalización de boletas de pago. 
Como metodología, se trabaja con un enfoque mixto que combina el análisis cualitativo mediante matrices FODA y entrevistas para el levantamiento de requisitos del negocio, 
junto con análisis cuantitativo de usabilidad, estructurado bajo una arquitectura web modular en capas (PHP/MySQL) con control transaccional atómico (COMMIT y ROLLBACK). 
Como resultado, la plataforma automatizó el cálculo masivo de haberes y la emisión instantánea de ocho tipos de reportes de dispersión y boletas digitales; asimismo, 
la evaluación empírica con la herramienta SUS (\textit{System Usability Scale}) arrojó una puntuación promedio de usabilidad del 87.8\% (calificación de excelencia, muy por encima del umbral funcional estándar de 68 puntos), 
registrando un 100\% de éxito en la integridad transaccional frente a fallos y errores de red. Finalmente, concluyen que la automatización de los procesos de cálculo e inclusión de mecanismos de preauditoría contable y control transaccional eliminan los descuadres financieros y reducen el tiempo operativo de las liquidaciones periódicas, permitiendo que la arquitectura modular y desacoplada incorpore nuevas normativas, facultades o conceptos remunerativos dinámicos sin necesidad de alterar el núcleo del código fuente.

% Articulo 2 RF 14

En este artículo de investigación publicado por \textcite{aliya2026design}, titulado ``Design and Pilot Evaluation of a Traceability-Oriented Enterprise Architecture for Advance Salary Payment Management'', donde mencionan que tiene como problemática la desconexión y fragmentación operativa entre los departamentos de recursos humanos, contabilidad y sistemas de pago externos en la gestión de anticipos salariales y haberes, donde la dependencia de aprobaciones en papel, cruces de correos y hojas de cálculo provoca duplicidad de registros, falta de sincronización con el sistema contable principal, retraso en la dispersión bancaria y una limitada trazabilidad de las transacciones para fines de auditoría. El objetivo es diseñar, implementar y evaluar una Arquitectura Empresarial Orientada a la Trazabilidad (TOEA) en tres capas que integre el autoservicio para empleados, la orquestación de flujos de aprobación administrativa, la sincronización contable automatizada mediante APIs y el monitoreo granular de eventos de pago. Como metodología, se utiliza la Metodología de Investigación en Ciencias del Diseño (DSR --- \textit{Design Science Research}) estructurada en seis fases de desarrollo y validación experimental, diseñando un prototipo funcional bajo una arquitectura en tres capas con identificación unificada mediante GUIDs y autenticación multifactor. Como resultado, se redujeron los pasos manuales del flujo de trabajo de 11 a 4 y se suprimieron en su totalidad las transferencias manuales de información entre áreas (de 6 a 0); además, en la validación con expertos se obtuvieron altas puntuaciones en auditabilidad y trazabilidad (4.81/5), transparencia operativa (4.75/5) y efectividad de sincronización (4.63/5), mientras que en pruebas de estrés con 1,000 peticiones concurrentes la plataforma mantuvo tiempos de respuesta estables de 1.97 segundos. Concluyen que la formalización de arquitecturas orientadas a la trazabilidad y la orquestación asíncrona de eventos resuelven de forma concluyente la fragmentación operativa de nómina, demostrando que el uso de identificadores GUIDs asociados a un registro inmutable garantiza la transparencia fiscal y financiera, ofreciendo un entorno escalable y extensible hacia analítica avanzada.

% Articulo 3 RF 15

En este artículo de investigación publicado por \textcite{tandrio2026web}, titulado ``Web-Based Payroll System Development Using the Prototyping Method and Structured Database Design'', donde mencionan que tiene como problemática la dependencia total operativa de hojas de cálculo en Excel para la elaboración de nómina y liquidación de incrementos salariales, condición semimanual que genera archivos corruptos, dispersión de datos entre finanzas y personal, consumo de 2 a 3 horas mensuales por parte de la dirección para liquidar haberes y una tasa de 0 a 6 errores de cálculo al mes, derivando en desbalances financieros, insatisfacción de los colaboradores y reclamos laborales. Tienen como objetivo desarrollar e implementar un sistema de información de nómina basado en web mediante el método de prototipado y un diseño estructurado de base de datos para centralizar la información salarial, automatizar el cálculo de pagos y erradicar errores aritméticos. Como metodología, se optó por la integración de la metodología de prototipado iterativo y el Ciclo de Vida del Desarrollo de Sistemas de Bases de Datos (DSDLC), iniciando con un diagnóstico situacional mediante el enfoque 5W+1H, complementado con diseño relacional normalizado y prototipado visual bajo el principio KISS. Como resultado, el diseño relacional estructurado requirió apenas 83.77 KB de almacenamiento mensual para el historial de transacciones, optimizando el rendimiento; asimismo, las pruebas de caja negra validaron el 100\% de los módulos sin fallas críticas y, en las pruebas UAT, el sistema obtuvo un índice de satisfacción global del 90.8\% (\textit{Highly Eligible}), destacando la claridad de la interfaz (4.6/5) y la exactitud en el cómputo automático de sueldos netos y préstamos frente al método manual en Excel. Concluyen que reemplazar el manejo de planillas en hojas de cálculo por una aplicación web con base de datos relacional normalizada suprime los márgenes de error de digitación y garantiza la integridad de los datos laborales, facilitando una adopción rápida entre el personal administrativo sin requerir entrenamientos complejos.

% Articulo 4 RF 16

En este artículo de investigación publicado por \textcite{someswari2025employee}, titulado ``Employee Payroll and Management System'', donde mencionan que tiene como problemática que la administración convencional de recursos humanos y nóminas se efectúa mediante procesos manuales o semiautomatizados que carecen de integración entre el registro de asistencia, políticas de la empresa y la remuneración salarial, lo cual ocasiona retrasos en los ciclos de aprobación, sobrecarga administrativa en el personal de RR.HH. por consultas repetitivas de los empleados, errores de cómputo en las deducciones y escasa transparencia en la entrega de comprobantes de pago. El principal objetivo es el diseño y despliegue de una plataforma web integral de gestión de empleados y planillas que automatice las tareas administrativas de personal, sincronice el control de asistencias con el cálculo salarial y proporcione un entorno interactivo y seguro de autoservicio para los trabajadores. Como metodología, se aplica una metodología de desarrollo iterativo para aplicaciones web, estructurada en una arquitectura desacoplada de tres niveles (presentación SPA, lógica de negocio y persistencia), incorporando control de acceso basado en roles (RBAC) y un asistente conversacional de políticas internas mediante arquitectura RAG. Como resultado, la automatización del cruce de asistencia con el motor de liquidación salarial y la emisión digital de boletas redujo el tiempo de procesamiento de nómina en aproximadamente un 50\%, eliminando tareas de digitación repetitiva; asimismo, el módulo de asistencias sostuvo índices de registro superiores al 90\%, y las funcionalidades de autoservicio del empleado permitieron alcanzar un nivel de satisfacción de usuario del 95\%. Concluyen que la centralización de las operaciones de RR.HH. y nómina en una plataforma web modular reduce sustancialmente el esfuerzo y las horas-hombre dedicadas al cálculo de las planillas, al tiempo que otorgar un portal de autoservicio seguro al empleado incrementa su nivel de satisfacción laboral y libera a las áreas contables de cargas administrativas rutinarias.

% Articulo 5 RF 17

En este artículo de investigación publicado por \textcite{chattopadhyay2026ai}, titulado ``AI-Enhanced Event-Driven Payroll Processing: 
A Cost-Optimized Framework Using Managed Container Services and Intelligent Automation'', 
donde menciona que tiene como problemática la ineficiencia económica y operativa en los sistemas tradicionales de nómina basados en servidores siempre encendidos 
(\textit{always-on}), los cuales presentan entre 97\% y 99\% de capacidad ociosa entre ciclos de pago quincenales o mensuales, 
sumado a la vulnerabilidad a fallas de cálculo salarial (tasas de error del 3\% al 8\% en validaciones manuales de asistencia y retenciones fiscales), 
altos costos de remediación y riesgo recurrente de sanciones tributarias por inconsistencias en los archivos de liquidación. 
Se propone experimentalmente un marco de trabajo de procesamiento de nómina orientado a eventos (\textit{event-driven}), 
optimizado en costos mediante servicios de contenedores administrados (Docker/MCS) y automatización inteligente de flujos con registro inmutable de auditoría. 
Como metodología, se aplica una metodología de diseño experimental y \textit{benchmarking} en la nube estructurada en cuatro fases operativas (ingesta, orquestación, procesamiento y entrega segura), 
validada experimentalmente mediante pruebas de rendimiento comparativo en entornos de contenedores sobre un corpus de 94 implementaciones empresariales y un estudio de caso longitudinal de 120 días. 
Como resultado, se redujo el costo de infraestructura de cómputo entre 71\% y 74\% frente al modelo tradicional de servidor dedicado; el tiempo medio de procesamiento de nómina (MTTP) 
se redujo de 6.4 horas a solo 38 minutos (90.1\% de reducción); la precisión en el cálculo salarial mejoró de 94.2\% a 99.7\%, logrando una disminución del 87\% en los errores operativos 
y reduciendo las horas de revisión manual de $\sim$420 horas a 54 horas por ciclo, a la par que las infracciones normativas cayeron de 2.3 a 0.04 por cada 100 ciclos (98.3\% de reducción). Finalmente, concluye que la sustitución de servidores dedicados por arquitecturas dirigidas mediante eventos y servicios de contenedores bajo demanda optimiza drásticamente los recursos computacionales y el gasto en infraestructura, garantizando que la orquestación asíncrona combinada con filtros de validación y pistas de auditoría inmutables asegura una exactitud salarial cercana al 100\% y neutraliza contingencias laborales y fiscales.